% ID: SEL-D0052-OA-E03
% ESTADO: validada. S053.
Sea $S$ el suceso «sigue sin fumar al cabo de un año». Las terapias A y B tienen probabilidades $0{,}45$ y $0{,}55$.

\textbf{a)} Por probabilidad total, $P(S)=0{,}45\cdot0{,}70+0{,}55\cdot0{,}80=\boxed{0{,}755}$.

\textbf{b)} Entre quienes siguen sin fumar, la fracción que hizo A es, por Bayes,
\[
P(A\mid S)=\frac{0{,}45\cdot0{,}70}{0{,}755}=\boxed{\tfrac{63}{151}}.
\]
\textbf{c)} En el caso de recaída, $P(S^c)=0{,}245$ y $P(S^c\mid A)=0{,}30$. Entonces
\[
P(A\mid S^c)=\frac{0{,}45\cdot0{,}30}{0{,}245}=\boxed{\tfrac{27}{49}}.
\]
